\originalchapter{线性规划建模与单纯形法}\label{chap:12}
本章补充 MIT 15.053 \emph{Optimization Methods in Management Science}, Spring 2013, James Orlin 与 Ebrahim Nasrabadi 课程中的第2--6、9讲及教程1、2、4--7. 它们补足6.253已讲的线性对偶理论所未展开的建模、标准形、单纯形计算、Phase I与灵敏度. 原课件中的休息插图、时事、项目安排不属于数学正文; 明示排除开放许可的出版社题面只保留来源指引, 相关数学方法在此独立推导.

\originalsection{变量、参数与线性模型}
模型要分别说明能控制的变量、预先给定的参数、比较方案的目标及限制方案的约束. 变量名称不是小事: 若约束跨多天、多个阶段或多个产品, 只记录单日总量可能丢失原问题的逻辑. 单位也应匹配; 时间必须统一为分钟或小时, 比例与总量不能随意相加. 线性模型采用可分性、比例性与确定系数等假设, 输出最优解后还要检查这些假设是否合理.

\begin{example}
教程1第5--9页的生产例: 产品P、Q分别需要机器A的50、24分钟和机器B的30、33分钟, 两台机器可用40、35小时; 单位利润为25、30美元. 写出连续生产模型.
\end{example}
\begin{solution}
令 $x,y$ 为两产品数量, 时间统一成分钟, 得 $\max(25x+30y)$, 约束 $50x+24y\le2400$, $30x+33y\le2100$, $x,y\ge0$. 原教程第7页的一行目标写成 $40x+35y$, 与同页说明及完整模型不符; 此处按单位利润25、30及后续完整模型计算. 若产品不可拆分, 应另要求整数; 连续模型只是其线性松弛, 不能把小数产量解释为可直接实施的整数方案.
\end{solution}

\begin{example}
教程1第10--17页的三产品模型, 机器时间矩阵与销售上限为
\[
A=\begin{pmatrix}20&10&10\\12&28&16\\15&6&16\\10&15&0\end{pmatrix},\quad
b=2400\mathbf1,\quad u=(100,40,60)\T.
\]
单位销售收入为 $(90,100,70)$, 材料成本为 $(45,40,20)$, 固定周费用6000美元. 求并验证最优连续生产方案.
\end{example}
\begin{solution}
模型为 $\max(45p+60q+50r-6000)$, $A(p,q,r)\T\le b$, $0\le(p,q,r)\T\le u$. 固定费用不影响最优变量, 但报告实际利润时不可遗漏. 取 $r=60$ 并使前两台机器满负荷, 解 $20p+10q=1800$, $12p+28q=1440$, 得 $p=900/11$, $q=180/11$. 后两台机器使用 $25140/11$ 与 $11700/11$ 分钟, 都不超过2400; 变量上限也满足.

为证明最优, 对前两台机器约束分别用乘子 $27/22,75/44$, 对 $r\le60$ 用乘子 $115/11$, 其余置零. 加权后各变量目标系数恰为45、60、50, 给上界 $84300/11$. 候选方案达到该界, 所以最优. 扣除6000后利润为 $18300/11$. 原教程的数值约7664及1664是四舍五入; 产品Q表中16.63与公式不一致, 正确近似为16.36.
\end{solution}

广告、营养、混合配料模型的共同结构是 $\min\sum_jc_jx_j$, $\ell_i\le\sum_ja_{ij}x_j\le u_i$, $0\le x_j\le d_j$. 索引集、系数与单位要在公式之前定义. 广告模型中 $a_j$ 表示一次展示带来的计数, 不自动等于不重复触达人数; 若同一受众会重复看到广告, 需说明计数假设. 混合配料中若总质量固定为 $P$, 成分约束可写 $L_jP\le\sum_i a_{ji}x_i+M_j\le U_jP$, 重量守恒为 $\sum_i x_i+\sum_jM_j=P$. 这里只将纯成分 $M_j$ 加在第 $j$ 个成分中, 不把所有纯材料同时加进每个成分约束. 教程2的外部熔炉书摘以这一抽象模型替代逐字重译, 原书出处保留在来源台账中.

\begin{example}
第2讲的历史饮食模型允许分数份量, 成本向量 $(1,3,2.5,3,1)$, 热量 $(250,770,360,190,230)$, 脂肪热量 $(81,360,144,45,99)$, 蛋白质 $(31,44,14,27,3)$, 钠 $(480,1170,800,580,160)$. 要求总热量600--900, 脂肪热量不超过总热量的40\%, 蛋白质至少30, 钠不超过1150. 写出模型并核对课件给的近似方案.
\end{example}
\begin{solution}
令份量向量 $x\ge0$, 总热量 $F=a\T x$. 目标 $\min c\T x$, 约束 $600\le F\le900$, $d\T x\le0.4F$, $p\T x\ge30$, $s\T x\le1150$. 仅用前两项, 令热量下界与脂肪比例取等号, 解得 $x_1=1040/921$, $x_2=380/921$, 其他为0, 成本 $2180/921\approx2.367$. 蛋白质及钠限制均满足. 可用热量与脂肪两条活跃约束构造非负对偶乘子证明其最优: 设热量下界乘子 $\alpha=109/27630$, 脂肪限制乘子 $\beta=2/2763$, 则 $250\alpha+19\beta=1$, $770\alpha-52\beta=3$; 其他食品的同一加权系数不超过其成本, 给下界 $600\alpha=2180/921$.

原第14页钠数据是1170, 第15页公式却写1770; 后者也与给定近似解不相容, 本稿按第14页数据建模. 此处数字是2013课程例题, 不作为现行营养、价格或饮食建议. 整数模型中一份第1项与两份第5项成本为3, 满足所有条件; 连续最优值2.367本身只给整数下界. 但成本低于3时, 第2与第4项不能选, 第3项最多一份且不能与任何其他正份量组合, 热量仅360; 或只选第1与第5项且总份数至多2, 热量最多500. 这些都违反600热量下界, 所以成本3确为整数最优.
\end{solution}

\begin{example}
第2讲的投资组合例给历史期望收益率 $(12.7,9.9,11.8)$ 与协方差矩阵
$Q=\left(\begin{smallmatrix}350&50&100\\50&150&30\\100&30&600\end{smallmatrix}\right)$. 写出收益--风险模型并计算权重 $(0.5,0.2,0.3)$ 的两个指标.
\end{example}
\begin{solution}
权重满足 $\mathbf1\T x=1,x\ge0$, 收益为 $r\T x$, 方差为 $x\T Qx$. 给定最大方差 $v$, 可以最大化 $r\T x$ 并要求 $x\T Qx\le v$. 协方差矩阵半正定使这是凸问题, 虽非线性规划. 样本权重收益为11.87, 方差为191.1. 每一对不同资产的协方差在二次型中出现两次, 不能只计算对角项. 改变风险上界形成有效前沿; 这是历史数学模型, 不涉及对实际证券作投资推荐.
\end{solution}

\originalsection{排班与辅助变量}
\begin{example}
第2讲与教程1的邮局例要求七天人数至少为 $(17,13,15,19,14,16,11)$, 每人连续工作五天再休息两天. 建立最少雇员模型并比较连续与整数方案.
\end{example}
\begin{solution}
令 $x_j$ 表示在第 $j$ 天开始五天班次的人数, 下标按模7循环. 第 $i$ 天上班人数为 $\sum_{k=0}^4x_{i-k}$, 因而模型是 $\min\sum_jx_j$, 约束 $\sum_{k=0}^4x_{i-k}\ge b_i$, $x\ge0$. 单纯定义每天上班总人数不能保证同一个人连续工作五天, 这就是选择班次变量的作用.

连续方案 $(4/3,10/3,2,22/3,0,10/3,5)$ 满足所有约束, 总数 $67/3$. 取非负乘子 $y=(1/3,0,1/3,1/3,0,1/3,0)$. 每个班次的覆盖乘子和都不超过1, 而需求加权和为 $67/3$. 因而对所有可行 $x$, $\sum_jx_j\ge y\T Ax\ge y\T b=67/3$, 候选方案达到下界, 证明连续最优.

整数可行方案 $(4,4,2,6,0,4,3)$ 总数23. 由连续最优值 $67/3$, 整数最优值至少为 $\lceil67/3\rceil=23$, 该方案遂最优. 原教程第23页写的向量把第六项写为6, 总和实际为25; 其后23周轮换表却给第六类4周. 本稿采用可行的4, 并给出上述连续下界证书以证明整数方案最优.

公平轮换可用长度23的班次序列, 各开始日重复次数依次为4、4、2、6、0、4、3. 让23名员工取该序列的23个循环平移, 每周各类人数保持不变, 每人在23周周期内经历相同班次多重集. 这证明了轮换的公平计数, 不自动处理跨周劳动法规或其他现实排班约束.
\end{solution}

若允许不足与超员, 可令每天 $w_i+d_i-e_i=b_i$, $d_i,e_i\ge0$, 并在目标加入正系数的不足与超员惩罚, 如 $5d_i^2+e_i^2$. 可行集中会出现 $d_i,e_i$ 同时正的“多余表示”, 但任何这样的解都可把两者同时减去 $\min(d_i,e_i)$, 保持等式而严格降低目标. 所以每个最优解至少一者为0, 正好具有原意. 若惩罚系数不正, 或目标不随这些变量严格增加, 这个等价性证明可能失效.

\originalsection{线性变换与标准形}
本章标准形统一为 $\max c\T x$, $Ax=b$, $x\ge0$, 并预处理使 $b\ge0$ 及 $A$ 行满秩. 极大与极小可通过目标换号相互转换. $a\T x\le b$ 加松弛变量 $s\ge0$ 变成 $a\T x+s=b$; $a\T x\ge b$ 减剩余变量 $s\ge0$ 变成 $a\T x-s=b$. 右端为负时先将整行乘以 $-1$, 同时反转不等号. $x_j\le0$ 可替换为 $-y_j$, $y_j\ge0$; 自由变量可拆成 $x_j=u_j-v_j$, $u_j,v_j\ge0$. 这一表示不唯一, 但投影回原变量时可行性与目标值完全相同.

\begin{example}
教程6的练习为 $\min(x_1-x_2+x_3)$, 约束 $x_1+2x_2-x_3\le3$, $-x_1+x_2+x_3\ge2$, $x_1-x_2=10$, $x_1\ge0,x_2\le0$, $x_3$ 自由. 转为极大标准形.
\end{example}
\begin{solution}
令 $y_2=-x_2$, $x_3=y_3-y_4$, 再加松弛及剩余变量, 得
\begin{align*}
\max\quad&-x_1-y_2-y_3+y_4,\\
 x_1-2y_2-y_3+y_4+s_1&=3,\\
 -x_1-y_2+y_3-y_4-s_2&=2,\\
 x_1+y_2&=10,
\end{align*}
所有新变量非负. 由等式消去自由变量也是可行方法, 但消去有符号限制的变量时必须保留对应不等式; 否则会把不可行方案引入. 目标常数即使不影响最优变量, 也要在恢复原目标值时计入.
\end{solution}

教程4还介绍三类常用变换. $|a\T x+b|\le d$ 等价于两条线性不等式, 要求 $d\ge0$; $|a\T x+b|\ge d>0$ 通常是两个半空间的并, 非凸, 不能用仅含连续变量的线性规划提升来一般表示, 因为线性像与投影保持凸性. 最小化有限个仿射函数的最大值可引入 $t\ge a_i\T x+b_i$ 并最小化 $t$; 最大化有限个仿射函数的最小值则令 $t\le a_i\T x+b_i$ 并最大化 $t$. 若反转优化方向, 新变量可能不被压到正确边界, 不能机械套用.

比例约束 $(a\T x+b)/(d\T x+e)\ge\rho$ 只有在分母已知严格正时才等价于 $a\T x+b\ge\rho(d\T x+e)$. 零分母处原比例未定义, 交叉相乘后的线性式却可能允许该点; 若采用这种扩展必须明确约定, 不能称严格等价. 教程中的“至少20\%广告属于某类”因此还要有正总广告量假设或明示零投放时的比例约定.

\originalsection{基本可行解与顶点}
设 $A\in\R^{m\times n}$ 行满秩. 选择 $m$ 个线性无关列组成可逆矩阵 $B$, 对应变量为基本变量, 其他为非基本变量. 设置 $x_N=0$ 后, $x_B=B^{-1}b$ 称基本解; 若它非负, 称基本可行解. 此时字典为
\[
x_B=B^{-1}b-B^{-1}A_Nx_N,\qquad
z=c_B\T B^{-1}b+\bar c_N\T x_N,
\]
其中 $\bar c_j=c_j-c_B\T B^{-1}a_j$ 是检验数或约化成本. 原课程的 $-z$ 行与这里的目标字典是同一种表达, 只需按等式解回 $z$, 不可直接凭行首符号猜方向.

\begin{theorem}
标准形可行点为顶点当且仅当其正坐标对应的列线性无关. 每个非空标准形可行集有基本可行解; 若最优值有限, 存在一个最优基本可行解.
\end{theorem}
\begin{proof}
若正坐标的列相关, 存在非零 $d$ 支持于这些坐标且 $Ad=0$. 足够小的正负步长保持非负, 点是两个不同可行点的中点, 不是顶点. 若列无关, 将 $x$ 写成可行点 $y,z$ 的中点, 非负性强制所有零坐标在 $y,z$ 也为0, 列无关又使 $A(y-z)=0$ 只能 $y=z$, 所以是顶点. 正坐标的独立列可补足为一个基, 因而顶点就是某个基本可行解.

任一可行点若支持列相关, 选择上述方向, 沿一个方向移动直到一个正坐标变为0, 得支持更小的可行点. 有限次后列独立. 若从最优点出发, 足够小的正负移动都可行, 最优性强制 $c\T d=0$, 因而沿方向降支持时保持最优. 有限最优值的取到已由多面体上的线性函数存在性证明, 所以得到最优基本可行解.
\end{proof}
一般多面体若含直线, 可以没有顶点. 若不含直线, 通过有限面下降或标准形的支持消去可得顶点; 在有有限最优值时, 最优面也不含直线, 因而有最优顶点. “顶点处必有最优解”需要非空、有限最优值及顶点存在, 不能用于不可行或无界问题. 开约束 $0<x<1$ 的 $\sup x=1$ 未取到, 也不属于闭多面体理论.

\originalsection{换基、检验数与终止证书}
当前基本可行解满足 $\bar c_j\le0$ 对全部非基本变量时, 任何可行解都有 $z=z_0+\sum_j\bar c_jx_j\le z_0$, 因而当前解最优. 这也给对偶证书 $\pi=B^{-\mathsf T}c_B$, 因为 $A\T\pi\ge c$, $b\T\pi=z_0$. 有正检验数时, 选择进入变量 $x_s$ 并增加至 $t\ge0$. 基本变量变为 $x_B(t)=\beta-t d$, $\beta=B^{-1}b$, $d=B^{-1}a_s$.

若所有 $d_i\le0$, 任意 $t\ge0$ 都可行而目标 $z_0+t\bar c_s\to\infty$, 给无界证书. 否则最大可行步长为 $t^*=\min_{d_i>0}\beta_i/d_i$, 达到该最小值的基本变量离基. 换基等价于以其行、进入列元素为主元作 Gauss--Jordan 消元: 主元行除以主元, 再消去其他行及目标行的同列系数. 只在正的 $d_i$ 上作比值检验, 负比值不能当作更小可行步长.

\begin{example}
第4讲用方程组 $2x_1+2x_2+x_3=9$, $2x_1-x_2+2x_3=6$, $x_1-x_2+2x_3=5$ 说明消元. 第5讲的字典为
\begin{align*}
z&=11-2x_2+6x_5,\\
x_3&=4-2x_2-2x_5,\quad x_4=1+x_2+2x_5,\\
x_1&=9-6x_2-3x_5.
\end{align*}
分别完成线性消元与单纯形的一次换基.
\end{example}
\begin{solution}
方程组先以首行的一半消去第一列, 得 $x_1+x_2+x_3/2=9/2$, $-3x_2+x_3=-3$, $-2x_2+3x_3/2=1/2$. 第二行规范化后消去第二列, 得 $x_1+5x_3/6=7/2$, $x_2-x_3/3=1$, $5x_3/6=5/2$. 最后 $x_3=3,x_2=2,x_1=1$.

字典中令 $x_5$ 进入. 正限制来自 $x_3$ 与 $x_1$, 比值为 $4/2=2$ 与 $9/3=3$, 所以 $x_3$ 离基. 解出 $x_5=2-x_2-x_3/2$, 代回得 $x_4=5-x_2-x_3$, $x_1=3-3x_2+3x_3/2$, $z=23-8x_2-3x_3$. 检验数都非正, 所以 $x_1=3,x_4=5,x_5=2$、其余0为最优, 值23. 若原字典改为 $z=6-2x_2+x_5$, $x_3=4-2x_2+2x_5$, $x_4=2+x_2+2x_5$, $x_1=3-6x_2$, 则令 $x_2=0,x_5\to\infty$ 给可行无界射线, 正好对应第5讲的无界例.
\end{solution}

\originalsection{Phase I与人工变量}
松弛变量表示原约束的剩余资源, 在最终解中可以为正. 人工变量只是为了构造初始基, 若它在辅助最优解中仍为正, 就不能当作原问题变量. 对 $Ax=b$, $x\ge0$, 先使 $b\ge0$, 加人工变量 $y\ge0$, 解
$\min\mathbf1\T y$, $Ax+y=b$, $x,y\ge0$. 初始基为 $y=b$. 目标非负, 且原问题可行当且仅当辅助最优值为0. 若原问题可行, 把其解与 $y=0$ 代入得0; 若辅助取到0, 非负性强制所有 $y=0$, 恢复原可行解.

若辅助值0但人工变量仍在基中且值为0, 可在该行选择一个非人工的非零系数换出它; 若没有这种系数, 该行是冗余等式, 可以在保留等价可行集的前提下移除. 然后删除人工列, 恢复原目标, 用基本行消去目标中的基本变量, 开始 Phase II. 原讲义把人工变量比作暂时的松弛, 但两者的语义与最终允许状态完全不同.

\begin{example}
第5讲的初始等式为 $x_1+x_2+x_3=4$, $-2x_1+x_2-x_3=1$, $x\ge0$. 构造并求解 Phase I.
\end{example}
\begin{solution}
加 $y_1,y_2$ 后初始值为 $(4,1)$, 辅助目标5. 首先让 $x_2$ 进入, 增至1使 $y_2=0$, 此时 $y_1=3$. 解出 $x_2=1+2x_1+x_3-y_2$, 代入另一行, 得 $y_1=3-3x_1-2x_3+y_2$, 目标为 $3-3x_1-2x_3+2y_2$. 令 $x_1$ 增至1, $y_1$ 变为0, $x_2=3,x_3=0$, 辅助值0. 所以原问题可行, 可使用 $x_1,x_2$ 作为起始基, 再按原目标继续.
\end{solution}

\originalsection{退化、循环与 Bland 规则}
基本可行解的某个基本变量为0时称退化. 最小比值可以为0, 此时换基但变量点和目标值不变. 同一个点可以对应多个基, 因而“每次换基就到了新点”不成立. 没有退化时每次正检验数换基都严格提高目标, 基只有有限个, 因而算法有限. 退化时需防循环.

\begin{example}
教程7第10--12页的循环例, 初始基为 $x_5,x_6,x_7$, 目标字典及三行约束为
\[
z=3+\tfrac34x_1-20x_2+\tfrac12x_3-6x_4,
\quad
\begin{pmatrix}
1/4&-8&-1&9&1&0&0\\
1/2&-12&-1/2&3&0&1&0\\
0&0&1&0&0&0&1
\end{pmatrix}x=\begin{pmatrix}0\\0\\1\end{pmatrix}.
\]
取进入序列 $x_1,x_2,x_3,x_4,x_5,x_6$, 离基序列 $x_5,x_6,x_1,x_2,x_3,x_4$, 验证循环.
\end{example}
\begin{solution}
每次按对应行的正主元消元, 所有比值都取0. 基序列依次为
\[
\begin{aligned}
&(5,6,7)\to(1,6,7)\to(1,2,7)\to(3,2,7),\\
&(3,2,7)\to(3,4,7)\to(5,4,7)\to(5,6,7).
\end{aligned}
\]
 第六次后矩阵与目标行恢复初始值, 变量点始终为 $x_7=1$、其余0, 目标始终3. 这说明合法的进入选择加普通比值规则可以无限循环, 并未说明所有选择规则都会循环. 来源课件对该例的历史作者归属未在本稿另行确认, 这里只使用已核的数学数据.
\end{solution}

\begin{theorem}[Bland 防循环规则]
极大单纯形中, 选择具有正检验数的最小编号非基本变量进入; 最小比值并列时选择编号最小的基本变量离开. 此规则不会循环, 因而在有限次换基后给最优或无界证书.
\end{theorem}
\begin{proof}
若有循环, 目标沿整个循环非降又回原值, 所以每次都退化, 可行变量点不变. 称循环中曾改变基本身份的变量为变化变量, 取编号最大的变化变量 $x_t$. 选一次 $x_t$ 离基、$x_s$ 进入的换基; 有 $s<t$. 该次进入方向 $d$ 满足 $d_s=1$, $d_i=-a_{is}$ 对当时基本变量, 其他非基本分量0, 且 $d_t<0$. 所有变化变量在固定可行点的值都为0. 若某个当时基本的变化变量 $i<t$ 有 $a_{is}>0$, 它也具有零最小比值, 按离基规则应先于 $t$ 离开, 矛盾. 因而对全部变化变量 $i<t$, 都有 $d_i\ge0$.

另选循环中 $x_t$ 进入的一个字典. 按进入规则, 该字典中编号小于 $t$ 的非基本变量检验数都非正, 而 $\bar c_t>0$. 编号大于 $t$ 的非基本变量不是变化变量, 始终非基本, 所以它们在前一个方向 $d$ 中为0. 同一满足 $Ad=0$ 的代数方向的目标变化可在任意字典中计算. 退化时它不一定允许正的可行步长, 此处只使用零空间恒等式. 在前一个字典它等于进入变量的正检验数 $\bar c_s>0$; 在后一个字典却为 $\sum_{j\in N}\bar c_jd_j<0$, 因为 $t$ 项严格负, 较小编号项非正, 较大编号项为0. 矛盾. 基有限, 无循环就保证有限性.
\end{proof}
教程7第15页把并列离基写成“编号最小行, 不考虑变量编号”; 这不是上述标准 Bland 规则. 若行随换基不按当前基本变量编号重新排序, 不能用行位置替代变量编号. 本稿采用可证明防循环的变量编号版本.

退化与多最优解也不同. 非基本变量有零检验数且允许正步长时, 可以沿该方向得到不同的同值可行解; 若只有零步长, 只能得到同一点的另一基. 因而“存在零检验数”本身不能无条件证明有不同的最优变量点.

\originalsection{基不变范围与影子价格}
固定最优基 $B$, 右端扰动 $b+\Delta b$ 时, 基本解为 $B^{-1}b+B^{-1}\Delta b$. 检验数不变, 因而只要这一向量仍非负, 原基仍最优, 最优值精确变化为 $\pi\T\Delta b$, 其中 $\pi=B^{-\mathsf T}c_B$. 对单个右端分量, 每个基本变量的非负条件给一个线性区间, 相交即有效范围. 超出范围后应换基, 旧影子价格可能仅为一个对偶界.

成本扰动 $c+\Delta c$ 不改变基本可行解, 但改变检验数. 新基仍最优当且仅当 $c_N+\Delta c_N-(c_B+\Delta c_B)\T B^{-1}A_N\le0$. 非基本变量单独改变利润时, 它的允许增加正好是负检验数的绝对值; 基本变量的改变会同时影响多个检验数, 需逐一求区间. 软件报告中的极大数如 $10^{30}$ 通常表示无上限, 不能把它当作可测的有限容许量.

\begin{example}
教程5的二维模型为 $\max(30K+50S)$, $2K+3S\le10$, $K+2S\le6$, $K+S\le5$, $K\le4$, $S\le3$, $K,S\ge0$. 求基点、第一条约束的影子价格与 $S$ 利润的允许变动范围.
\end{example}
\begin{solution}
前两条取等号得 $(K,S)=(2,2)$, 值160. 乘子 $(10,10)$ 对前两约束给目标法向量 $(30,50)$, 其他置零, 所以由对偶证书知最优. 将第一右端改为 $10+\Delta$, 同一交点为 $K=2+2\Delta$, $S=2-\Delta$, 值 $160+10\Delta$. 所有其余不等式及非负性给 $-1\le\Delta\le1$, 故影子价格10在该范围精确有效.

若目标变为 $30K+(50+\eta)S$, 保持该点最优要求 $(30,50+\eta)$ 是活跃法向量 $(2,3),(1,2)$ 的非负组合. 解乘子得第一为 $10-\eta$, 第二为 $10+2\eta$, 因而 $-5\le\eta\le10$. 两端某个乘子为0, 目标与边平行, 可出现多最优解. 区间内部基点不变而目标值按 $2\eta$ 改变.
\end{solution}

若多个右端同时改变, 单项允许范围不能逐项独立套用. 设方向对应的单项最大幅度为 $r_i>0$, 若 $\sum_i|\Delta b_i|/r_i\le1$, 则扰动向量是零扰动与各轴允许端点的凸组合. 同一基可行右端区域由线性不等式定义而凸, 所以新右端仍可行, 检验数不变, 影子价格仍有效. 这是100\%规则的证明; 它是充分条件, 不是必要条件. 无穷允许幅度可通过有限近似端点取极限处理.

\begin{example}
第6讲虚构的 mc2 生产例最优变量为 $(18,16.4,3,0,32)$, 利润46.12百万美元. 其新磁盘约束影子价格0.8、右端20、允许减少16.4及增加0.6; 产品A需求上限的影子价格1.04、右端18、允许减少18及增加0.75; 产品C上限影子价格0.6、右端3、允许减少3及增加0.6. 应用这些报告数字.
\end{example}
\begin{solution}
单位产量为千件, 利润单位为百万美元. 新磁盘从20降至15仍在 $[3.6,20.6]$ 内, 利润变化为 $-5\times0.8=-4$ 百万美元. 将A需求上限从18增至18.5, 变化 $0.5\times1.04=0.52$ 百万美元; 若广告成本0.4百万美元, 在其余假设不变时净增加0.12百万美元. 这只是模型内比较, 不预设现实需求一定按假设改变.

同时把A上限降至15、C上限降至2, 100\%比率为 $3/18+1/3=1/2$, 所以原影子价格有效, 利润降低 $3\times1.04+1\times0.6=3.72$ 百万美元. 若A降1而C增0.5, 比率 $1/18+0.5/0.6<1$, 利润变化 $-1.04+0.3=-0.74$ 百万美元. 非活跃资源约束的影子价格为0, 小范围增加该资源并不提高利润; 活跃约束也可能乘子为0, 不能把“活跃”与“有价值”当作等价.
\end{solution}

约化成本还有经济解释 $\bar c_j=c_j-\pi\T a_j$: 收入减去按影子价格计的资源成本. mc2 的产品D收益0.6, 使用0.5单位新磁盘及1单位电子阅读器容量, 对应资源价0.8及0.3, 所以 $\bar c_D=0.6-0.5\times0.8-0.3=-0.1$. 收益升至0.7才消除负检验数; 等号处未必立即要求正生产, 还要看允许进入方向与退化. 新产品F收益0.65、使用0.4单位新磁盘与1单位阅读器容量, 约化成本为 $0.65-0.4\times0.8-0.3=0.03>0$, 表示当前对偶价格不再能证明原方案最优, 可以通过换基寻求改进. 若模型同时改变多个成本或矩阵中的基列, 必须重新检查全部基可行性与检验数, 单一报告范围不足以说明结论.
